\subsubsection{Weak Convergence}\label{subsec:weak-convergence}

Procedures to solve differential equations usually involve a sequence of solutions to approximate problems. Very often it is possible to obtain uniform bounds for the approximations in the norm of some Banach space, but it is not so easy to show that they actually converge in that Banach space. In these situations, a weaker notion of convergence, known as weak convergence, is extremely useful. Weak convergence of a sequence of functions means that certain linear functionals, when applied to that sequence, yield convergent sequences of numbers. We have already encountered such a notion when we defined a sequence of distributions $f_n$ to be convergent if $(f_n,\phi)$ converges for every test function $\phi$. In the context of Banach spaces, we make the following definition.

\begin{Definition}
\label{def:weak-and-weak-star-convergence}
Let $X$ be a Banach space. A sequence $x_n$ in $X$ \textbf{converges weakly} to $x$ if $f(x_n)$ converges to $f(x)$ for every $f\in X^*$. A sequence $f_n$ in $X^*$ \textbf{converges weakly-*} to $f$ if $f_n(x)$ converges to $f(x)$ for every $x\in X$.
\end{Definition}

To distinguish notations, one writes $x_n\to x$ for convergence in norm, $x_n\rightharpoonup x$ for weak convergence, and $x_n\overset{*}{\rightharpoonup}x$ for weak-* convergence.

\begin{Remark}
\label{rem:weak-convergence-consequences}
We list here a number of important consequences of the definition of weak convergence. We give only very brief indications of the (usually very short) proofs of these assertions.\begin{itemize}
\item It follows directly from the definition that a weakly convergent sequence in $X^*$ also converges weakly-*. The converse is false in general, but it is true if $X$ is reflexive.
\item The uniqueness of weak limits follows from the Hahn-Banach theorem; the uniqueness of weak-* limits is obvious from the definition.
\item It follows from the uniform boundedness principle that weakly or weakly-* convergent sequences are bounded.
\item If $f(x_n)$ is a Cauchy sequence for every $f\in X^*$, then
\[l(f):=\lim_{n\to\infty}f(x_n)\]
defines a linear functional on $X^*$, but we do not necessarily know that this linear functional can be represented in the form $l(f)=f(x)$, where $x\in X$. Thus a Cauchy criterion does not necessarily hold for weak convergence. The spaces for which it does hold are called weakly (sequentially) complete. Obviously, reflexive spaces are weakly complete. On the other hand, a Cauchy criterion always holds for weak-* convergence.
\end{itemize}
\end{Remark}

\Needspace{4\baselineskip}
\begin{Example}
\label{ex:orthonormal-basis-weak-zero-example}
Let $H$ be a separable, infinite-dimensional Hilbert space and let $e_n$ be an orthonormal basis of $H$. Then $e_n$ converges weakly to zero (Problem~\ref{ex:orthonormal-basis-weak-zero}).
\end{Example}

\Needspace{4\baselineskip}
\begin{Example}
\label{ex:oscillatory-sines-weak-zero}
A particular case of the previous example is the sequence $f_n(x)=\sin nx$ in $L^2(0,1)$. This sequence of increasingly oscillatory functions converges weakly to zero. Of course, it does not converge to zero in any pointwise sense, though one could say it converges to zero ``on average.''
\end{Example}
